% Generated draft: requires mathematical and source review.
\documentclass[11pt]{article}
\usepackage[margin=25mm]{geometry}
\usepackage{fontspec}
\setmainfont{Latin Modern Roman}
\usepackage{amsmath,amssymb,mathtools,unicode-math}
\setmathfont{Latin Modern Math}
\usepackage{enumitem,needspace,xurl,hyperref}
\setlength{\parindent}{0pt}
\setlength{\parskip}{5pt}
\setlength{\emergencystretch}{3em}
\newcommand{\N}{\mathbb{N}}
\newcommand{\Z}{\mathbb{Z}}
\newcommand{\Q}{\mathbb{Q}}
\newcommand{\R}{\mathbb{R}}
\newcommand{\C}{\mathbb{C}}
\newcommand{\F}{\mathbb{F}}
\newcommand{\A}{\mathbb{A}}
\newcommand{\PP}{\mathbb{P}}
\newcommand{\E}{\mathbb{E}}
\newcommand{\eps}{\varepsilon}
\newcommand{\dd}{\,\mathrm{d}}
\begin{document}
\section*{UnsolvedMath ID: 30002393}
\textbf{Source problem code:} catalogue entry
\subsection{Definitions and mathematical statement}
For normed algebras \(A\) and \(B\) over a common scalar field, a linear map \(T\colon A\to B\) is bounded if
\[
\|T\|:=\sup_{\substack{a\in A\\ \|a\|\leq 1}}\|T(a)\|<\infty.
\]
Equivalently, there exists \(C\geq 0\) such that \(\|T(a)\|\leq C\|a\|\) for every \(a\in A\).

The supplied source asks: What methods produce bounded linear maps between close, or completely close, operator algebras?

This is a methodological question with unresolved formulation details. The supplied record does not specify the class of operator algebras, their representations, the meaning or quantitative threshold of either closeness condition, or the additional properties required of the constructed maps. In particular, the zero map \(T\colon A\to B\), defined by \(T(a)=0\) for every \(a\in A\), is always bounded and linear. Thus existence of a bounded linear map alone does not express a substantive perturbation problem. Any intended nontriviality, approximation, invertibility, multiplicativity, or quantitative norm requirement must be recovered from the source context before a precise theorem-like target can be stated. No such requirement is imposed in this draft.
\subsection{Short English statement}
The source asks for methods constructing bounded linear maps between operator algebras that are close or completely close. It leaves the closeness conditions and the desired properties of the maps unspecified; these are essential because the zero map already satisfies boundedness and linearity.
\subsection{Sources}
\begin{itemize}
\item Supplied local UnsolvedMath registry record, permanent ID \(30002393\), source problem code \texttt{catalogue entry}, title: Bounded Maps Between Close Operator Algebras. The record lists proposed year 2013 and identifies no individual proposer.
\item Supplied report citation: \emph{C*-Algebren (2014)}, Oberwolfach Reports, pp. 2423--2500, DOI \texttt{10.4171/owr/2013/43}. Original OWR report: \url{https://doi.org/10.4171/owr/2013/43}. The supplied background describes Oberwolfach Reports as documenting mathematical workshops and their open-problem discussions.
\item Supplied literature reference, labeled ``Christensen et al. primary paper'': \url{https://doi.org/10.1007/s11511-012-0075-5}. No paper title, full author list, or specific result is supplied.
\item Literature assessment dated exactly \texttt{2026-08-21}: the supplied curation labels the problem open, describes bounded-map constructions as known in special perturbation settings, and describes the general methodological problem as remaining open. This is supplied dated curation, not an independently verified assessment. The record explicitly requests independent verification of time-sensitive status and cited claims.
\item The registry labels the statement exact and reports that its mathematical content was checked against the official Oberwolfach report, with notation, context, and question form normalized. This verification claim is supplied provenance and has not been independently checked for this draft.
\end{itemize}
\end{document}
